% !TeX program = xelatex
% UTF-8; compile twice with XeLaTeX. No external images or data files needed.
\documentclass[a4paper,landscape,10pt]{article}
\usepackage[margin=14mm,headheight=14pt,headsep=5mm,footskip=7mm]{geometry}
\usepackage{fontspec}
\setmainfont{DejaVu Sans}
\newfontfamily\ipafont{DejaVu Sans}
\usepackage[table]{xcolor}
\usepackage{array,longtable,booktabs,fancyhdr}
\definecolor{ink}{HTML}{18384D}
\definecolor{band}{HTML}{E7EFF4}
\pagestyle{fancy}\fancyhf{}
\fancyhead[L]{\small\color{ink}Parlez Vous QSO?! · Català / Catalan}
\fancyhead[R]{\small Englische Phrase · Traducció catalana · IPA}
\fancyfoot[L]{\scriptsize Vorlage: Catalan v 1.0 · Credits: Jaume -- EA3FRB}
\fancyfoot[R]{\small\thepage}
\setlength{\parindent}{0pt}
\setlength{\tabcolsep}{3mm}
\renewcommand{\arraystretch}{1.25}
\newcommand{\ipa}[1]{{\ipafont[#1]}}
\begin{document}
{\LARGE\bfseries\color{ink}Parlez Vous QSO?!}\hfill{\large Català / Catalan}

\vspace{2mm}
{\small QSO-Sprachführer mit katalanischer IPA-Lautschrift. Die Lautschrift bezieht sich auf die katalanische Spalte.}

\vspace{2mm}
{\footnotesize \textbf{Leseschlüssel:} \ipa{ˈ} Hauptbetonung, \ipa{ˌ} Nebenbetonung; \ipa{ə} unbetonter Zentralvokal;
\ipa{ɛ, ɔ} offene e/o, \ipa{ʎ} ll, \ipa{ɲ} ny, \ipa{ʃ} ix/x, \ipa{ʒ} j/g vor e/i.
Grundlage: Zentralkatalanisch (Raum Barcelona). Schrägstriche trennen Alternativen; \ipa{…} steht für einen Platzhalter.}

\vspace{1mm}
\fontsize{9}{11.8}\selectfont
\begin{longtable}{@{}>{\raggedright\arraybackslash}p{13mm}>{\raggedright\arraybackslash}p{73mm}>{\raggedright\arraybackslash}p{75mm}>{\raggedright\arraybackslash}p{90mm}@{}}
\toprule
\textbf{Ref.} & \textbf{English phrase} & \textbf{Català} & \textbf{Lautschrift (IPA)}\\
\midrule
\endfirsthead
\toprule
\textbf{Ref.} & \textbf{English phrase} & \textbf{Català} & \textbf{Lautschrift (IPA)}\\
\midrule
\endhead
\midrule
\multicolumn{4}{r}{\scriptsize Fortsetzung auf der nächsten Seite}\\
\endfoot
\bottomrule
\endlastfoot
\rowcolor{band} 1.0 & \textbf{Opening Greetings} & \multicolumn{2}{l}{\textbf{Salutacions inicials}} \\*
1.1 & I am very pleased to meet you. & Estic content de conèixer-lo. & \ipa{əsˈtik kunˈtɛn də kuˈnɛʃəlu} \\
1.2 & Thank you for your call & Gràcies per la seva crida. & \ipa{ˈɡɾasiəs pəɾ lə ˈsebə ˈkɾidə} \\
1.3 & Good morning & Bon dia. & \ipa{ˈbɔn ˈdiə} \\
1.31 & Good day & Bon dia. & \ipa{ˈbɔn ˈdiə} \\
1.32 & Good afternoon & Bona tarda (12–18 h). & \ipa{ˈbɔnə ˈtaɾdə} \\
1.33 & Good evening & Bon vespre & \ipa{ˈbɔn ˈbɛspɾə} \\
1.4 & I know only a little Catalan but I enjoy QSOs in Catalan. & Sols parlo un xic de català però gaudeixo dels QSO en català. & \ipa{ˈsɔls ˈpaɾlu un ˈʃik də kətəˈla pəˈɾɔ ɡəwˈdɛʃu dəls ku ˈɛsə o ən kətəˈla} \\
\rowcolor{band} 2.0 & \textbf{Operator} & \multicolumn{2}{l}{\textbf{Operador}} \\*
2.1 & My name is Peter. & El meu nom és Pere. & \ipa{əl ˈmɛw ˈnɔm ˈes ˈpeɾə} \\
2.2 & I have been licensed since 1960. & Tinc la llicència des de l’any 1960. & \ipa{ˈtiŋ lə ʎiˈsɛnsiə ˈdez də ˈlaɲ ˈmil ˌnɔwˈsens ʃəˈʃantə} \\
2.3 & I have been licensed for 43 years. & Tinc la llicència de fa 43 anys. & \ipa{ˈtiŋ lə ʎiˈsɛnsiə də ˈfa kuˈɾantə ˈtɾɛs ˈaɲs} \\
2.4 & I live in Sittingbourne. & Visc a Sittingbourne. & \ipa{ˈbisk ə ˈsɪtɪŋbɔːn} \\
2.5 & Sittingbourne is a village. & Sittingbourne és un llogaret. & \ipa{ˈsɪtɪŋbɔːn ˈes un ʎuɡəˈɾɛt} \\
2.51 & Sittingbourne is a town. & Sittingbourne és un poble. & \ipa{ˈsɪtɪŋbɔːn ˈes un ˈpɔblə} \\
2.52 & Sittingbourne is a city. & Sittingbourne és una ciutat. & \ipa{ˈsɪtɪŋbɔːn ˈes ˈunə siwˈtat} \\
2.6 & Sittingbourne is between London and Dover. & Sittingbourne està entre Londres i Dover. & \ipa{ˈsɪtɪŋbɔːn əsˈta ˈɛntɾə ˈlondɾəs i ˈdəʊvə} \\
2.7 & I live in the South of England. & Visc al sud d’Anglaterra. & \ipa{ˈbisk əl ˈsut dəŋɡləˈtɛrə} \\
2.71 & I live in the North of England. & Visc al nord d’Anglaterra. & \ipa{ˈbisk əl ˈnɔɾt dəŋɡləˈtɛrə} \\
2.72 & I live in the East of England. & Visc a l’est d’Anglaterra. & \ipa{ˈbisk ə ˈlɛst dəŋɡləˈtɛrə} \\
2.73 & I live in West of England. & Visc a l’oest d’Anglaterra. & \ipa{ˈbisk ə luˈɛst dəŋɡləˈtɛrə} \\
2.74 & I live in South East England. & Visc al sud-est d’Anglaterra. & \ipa{ˈbisk əl ˌsudˈɛst dəŋɡləˈtɛrə} \\
2.75 & I live in South West England. & Visc al sud-oest d’Anglaterra. & \ipa{ˈbisk əl ˌsuduˈɛst dəŋɡləˈtɛrə} \\
2.76 & I live in North East England. & Visc al nord-est d’Anglaterra. & \ipa{ˈbisk əl ˌnɔɾdˈɛst dəŋɡləˈtɛrə} \\
2.77 & I live in North West England. & Visc al nord-oest d’Anglaterra. & \ipa{ˈbisk əl ˌnɔɾduˈɛst dəŋɡləˈtɛrə} \\
2.78 & I live in central England. & Visc al centre d’Anglaterra. & \ipa{ˈbisk əl ˈsɛntɾə dəŋɡləˈtɛrə} \\
2.8 & I live in the country. & Visc a l’interior. & \ipa{ˈbisk ə lintəɾiˈo} \\
2.81 & I live by the sea. & Visc vora el mar. & \ipa{ˈbisk ˈbɔɾə əl ˈmaɾ} \\
2.82 & I live near mountains. & Visc a muntanya. & \ipa{ˈbisk ə munˈtaɲə} \\
2.9 & Sittingbourne is 50Km from London. & Sittingbourne és a 50 Km. de Londres. & \ipa{ˈsɪtɪŋbɔːn ˈes ə siŋˈkwantə kiluˈmɛtɾəs də ˈlondɾəs} \\
2.11 & I am retired. & Estic jubilat. & \ipa{əsˈtik ʒubiˈlat} \\
2.12 & My home call sign is M0ECS & El meu indicatiu és M0ECS. & \ipa{əl ˈmɛw indikəˈtiw ˈes ˈɛmə ˈzɛɾu e se ˈɛsə} \\
\rowcolor{band} 3.0 & \textbf{Station} & \multicolumn{2}{l}{\textbf{Estació}} \\*
3.1 & My radio is an Icom 706. & La meva ràdio és una Icom 706. & \ipa{lə ˈmɛbə ˈradiu ˈes ˈunə ˈaɪkɒm ˌsɛtˈsens ˈsis} \\
3.2 & My station is all homebrew (homemade). & La meva estació és autoconstruïda. & \ipa{lə ˈmɛbə əstəsiˈo ˈes ˌawtukunstɾuˈidə} \\
3.3 & I am running 5 Watts. & Transmeto amb 5 Watt. & \ipa{tɾənzˈmɛtu əm ˈsiŋ ˈwats} \\
3.4 & My antenna is a long wire. & La meva antena és de fil llarg. & \ipa{lə ˈmɛbə ənˈtɛnə ˈes də ˈfil ˈʎaɾk} \\
3.41 & My antenna is a dipole. & La meva antena és un dipol. & \ipa{lə ˈmɛbə ənˈtɛnə ˈes un diˈpɔl} \\
3.42 & My antenna is a G5RV. & La meva antena és una G5RV. & \ipa{lə ˈmɛbə ənˈtɛnə ˈes ˈunə ʒe ˈsiŋ ˈɛrə ˈbe ˈbaʃə} \\
3.43 & My antenna is a doublet. & La meva antena és de doblet. & \ipa{lə ˈmɛbə ənˈtɛnə ˈes də duˈblɛt} \\
3.44 & My antenna is a delta loop. & La meva antena és una delta loop. & \ipa{lə ˈmɛbə ənˈtɛnə ˈes ˈunə ˈdɛltə luːp} \\
3.45 & My antenna is a vertical. & La meva antena és una vertical. & \ipa{lə ˈmɛbə ənˈtɛnə ˈes ˈunə bəɾtiˈkal} \\
3.46 & My antenna is a log periodic. & La meva antena és una periòdica logarítmica. & \ipa{lə ˈmɛbə ənˈtɛnə ˈes ˈunə pəɾiˈɔdikə luɡəˈɾitmikə} \\
3.47 & My antenna is a Yagi beam. & La meva antena és una Yagi directiva. & \ipa{lə ˈmɛbə ənˈtɛnə ˈes ˈunə ˈjaɡi diɾəkˈtibə} \\
3.48 & My antenna is a 5 element beam. & La meva antena és una directiva de 5 elements. & \ipa{lə ˈmɛbə ənˈtɛnə ˈes ˈunə diɾəkˈtibə də ˈsiŋ ələˈmens} \\
\rowcolor{band} 4.0 & \textbf{Conditions} & \multicolumn{2}{l}{\textbf{Condicions}} \\*
4.1 & Thank you for my 599 report. & Gràcies per l’informe 599. & \ipa{ˈɡɾasiəs pəɾ linˈfɔɾmə ˈsiŋ ˈnɔw ˈnɔw} \\
4.2 & The band conditions are good. & Les condicions a la banda són bones. & \ipa{ləs kundisiˈons ə lə ˈbandə ˈson ˈbɔnəs} \\
4.3 & The band conditions are bad. & Les condicions a la banda són dolentes. & \ipa{ləs kundisiˈons ə lə ˈbandə ˈson duˈlentəs} \\
4.4 & There is strong interference here. & Hi ha una interferència forta aquí. & \ipa{i ˈa ˈunə intəɾfəˈɾɛnsiə ˈfɔɾtə əˈki} \\
4.5 & Can you change to the 80m band? & Pot canviar a la banda de 80 m.? & \ipa{ˈpɔt kəmˈbja ə lə ˈbandə də bujˈtantə ˈmɛtɾəs} \\
4.6 & Please change frequency to 3.558 MHz. & Si us plau, canviï a la freqüència de 3.558 MHz. & \ipa{si us ˈplaw kəmˈbii ə lə fɾəˈkwɛnsiə də ˈtɾɛs ˈkɔmə ˈsiŋ ˈsiŋ ˈbujt ˈmɛɡəˌhɛɾts} \\
\rowcolor{band} 5.0 & \textbf{General Queries and Comments} & \multicolumn{2}{l}{\textbf{Preguntes i comentaris generals}} \\*
5.1 & What does <\#\#\#\#\#> mean please? & Què significa <\#\#\#\#\#>? & \ipa{ˈkɛ siɡniˈfikə …} \\
5.2 & Please wait a moment, I need to check my Catalan dictionary. & Esperi un moment si us plau. Necessito consultar el diccionari de català. & \ipa{əsˈpɛɾi un muˈmen si us ˈplaw nəsəˈsitu kunsulˈta əl diksiuˈnaɾi də kətəˈla} \\
5.3 & How do you say, <\#\#\#\#\#> in Catalan ? & Com es diu <\#\#\#\#\#> en català? & \ipa{ˈkɔm əs ˈdiw … ən kətəˈla} \\
5.5 & My favourite modes of operation are CW / RTTY / PSK / MT63 / Olivia / Hellschreiber / Digital Modes / and SSB & Les meves modalitats d’operació favorites són CW / RTTY / PSK / MT63 / Olivia / Hellschreiber / modes digitals / i SSB. & \ipa{ləs ˈmɛbəs mudəliˈtats dupəɾəsiˈo fəbuˈɾitəs ˈson se ˈbe ˈdɔblə / ˈɛrə te te i ˈɡɾɛɡə / pe ˈɛsə ˈka / ˈɛmə te ʃəˈʃantə ˈtɾɛs / uˈlibiə / ˈhɛlˌʃʁaɪbɐ / ˈmɔdəs diʒiˈtals / i ˈɛsə ˈɛsə be} \\
5.6 & My best DX is Australia. & El meu millor DX és Austràlia. & \ipa{əl ˈmɛw miˈʎo de ˈiks ˈes əwsˈtɾaliə} \\
5.7 & I enjoy working contests. & Frueixo participant en concursos. & \ipa{fɾuˈɛʃu pəɾtisiˈpan ən kuŋˈkuɾsus} \\
5.8 & I like working QRP. & M’agrada transmetre en QRP. & \ipa{məˈɡɾadə tɾənzˈmetɾə ən ku ˈɛrə pe} \\
\rowcolor{band} 6.0 & \textbf{The Weather} & \multicolumn{2}{l}{\textbf{El temps}} \\*
6.1 & It is sunny. & Dia assolellat. & \ipa{ˈdiə əsuʎəˈʎat} \\
6.2 & It is windy. & Dia ventós. & \ipa{ˈdiə bənˈtos} \\
6.3 & It is raining. & Plou. & \ipa{ˈplɔw} \\
6.4 & It is snowing. & Neva. & \ipa{ˈnɛbə} \\
6.5 & It is very hot. & Fa molta calor. & \ipa{ˈfa ˈmɔltə kəˈlo} \\
6.6 & It is very cold. & Fa molt fred. & \ipa{ˈfa ˈmol ˈfɾɛt} \\
6.7 & The temperature is 8C. & La temperatura és 8º C. & \ipa{lə təmpəɾəˈtuɾə ˈes ˈbujt ˈɡɾaws ˈsɛlsiws} \\
6.8 & It is overcast. & Està núvol. & \ipa{əsˈta ˈnubul} \\
6.9 & It is foggy. & Hi ha boira. & \ipa{i ˈa ˈbɔjɾə} \\
6.11 & It has been raining all day. & Ha estat plovent tot el dia. & \ipa{ə əsˈtat pluˈben ˈtot əl ˈdiə} \\
6.12 & There is 20cm of snow here. & Hi ha 20 cm. de neu aquí. & \ipa{i ˈa ˈbin sənˈtimətɾəs də ˈnɛw əˈki} \\
\rowcolor{band} 8.0 & \textbf{QSL} & \multicolumn{2}{l}{\textbf{QSL}} \\*
8.1 & My QSL card is sure via the bureau & La meva targeta QSL amb tota seguretat via bureau (associació). & \ipa{lə ˈmɛbə təɾˈʒɛtə ku ˈɛsə ˈɛlə əm ˈtotə səɡuɾəˈtat ˈbiə byˈʁo} \\
8.2 & QSL via eQSL & QSL via eQSL. & \ipa{ku ˈɛsə ˈɛlə ˈbiə i ku ˈɛsə ˈɛlə} \\
8.3 & No eQSL please & No eQSL, si us plau. & \ipa{ˈno i ku ˈɛsə ˈɛlə si us ˈplaw} \\
8.4 & I will send you my card direct. & Li enviaré la meva targeta directament. & \ipa{li əmbiəˈɾe lə ˈmɛbə təɾˈʒɛtə diˌɾɛktəˈmen} \\
8.5 & What is your address? & Quina és la seva adreça? & \ipa{ˈkinə ˈes lə ˈsebə əˈdɾɛsə} \\
8.6 & Please send your QSL card direct. & Si us plau, enviï la seva QSL directament. & \ipa{si us ˈplaw əmˈbii lə ˈsebə ku ˈɛsə ˈɛlə diˌɾɛktəˈmen} \\
8.7 & My address is... & La meva adreça és… & \ipa{lə ˈmɛbə əˈdɾɛsə ˈes …} \\
8.8 & My address is shown on QSL.com / EHam.net & La meva adreça està disponible a QSL.com / EHam.net. & \ipa{lə ˈmɛbə əˈdɾɛsə əsˈta dispuˈniblə ə ku ˈɛsə ˈɛlə ˈpun ˈkom / i hæm ˈpun ˈnɛt} \\
8.9 & Sorry I do not send QSL cards. & Ho sento. No envio targetes QSL. & \ipa{u ˈsɛntu ˈno əmˈbiu təɾˈʒɛtəs ku ˈɛsə ˈɛlə} \\
8.11 & Please send QSL card via G1JQH. & Si us plau, targeta QSL via G1JQH. & \ipa{si us ˈplaw təɾˈʒɛtə ku ˈɛsə ˈɛlə ˈbiə ʒe ˈun ˈʒɔtə ku ˈak} \\
8.12 & My WACRAL / GQRP / DXCC / FISTS number is 1234. & El meu número WACRAL / GQRP / DXCC / FISTS és 1234. & \ipa{əl ˈmɛw ˈnuməɾu ˈwakɾəl / ʒe ku ˈɛrə pe / de ˈiks se se / fɪsts ˈes ˈmil dusˈsens ˈtɾɛntə ˈkwatɾə} \\
\rowcolor{band} 9.0 & \textbf{Final Greetings} & \multicolumn{2}{l}{\textbf{Acomiadament}} \\*
9.1 & Thank you. & Gràcies. & \ipa{ˈɡɾasiəs} \\
9.2 & Thank you for a very enjoyable contact. & Gràcies per un contacte tan agradable. & \ipa{ˈɡɾasiəs pəɾ un kunˈtaktə ˈtan əɡɾəˈdablə} \\
9.3 & I enjoyed our chat. & He gaudit de la conversa. & \ipa{e ɡəwˈdit də lə kumˈbɛɾsə} \\
9.4 & I enjoyed our QSO in Catalan. & M’ha agradat el nostre QSO en català. & \ipa{mə əɡɾəˈdat əl ˈnɔstɾə ku ˈɛsə o ən kətəˈla} \\
9.5 & It has been great fun. & Ha sigut divertit. & \ipa{ə siˈɡut dibəɾˈtit} \\
9.6 & If you would like a regular contact with me, you can contact me via email. & Si li agradés un contacte regular amb mi, pot contactar-me via correu electrònic. & \ipa{si li əɡɾəˈdes un kunˈtaktə rəɡuˈla əm ˈmi ˈpɔt kuntəkˈtaɾmə ˈbiə kuˈrɛw ələkˈtɾɔnik} \\
9.7 & My email address is... & La meva adreça de correu electrònic és… & \ipa{lə ˈmɛbə əˈdɾɛsə də kuˈrɛw ələkˈtɾɔnik ˈes …} \\
9.8 & My email address is shown on QSL.com / EHam.net & La meva adreça de correu electrònic està disponible a QSL.com / EHam.net. & \ipa{lə ˈmɛbə əˈdɾɛsə də kuˈrɛw ələkˈtɾɔnik əsˈta dispuˈniblə ə ku ˈɛsə ˈɛlə ˈpun ˈkom / i hæm ˈpun ˈnɛt} \\
9.9 & I look forward to meeting you again. & Espero retrobar-te novament. & \ipa{əsˈpɛɾu rətɾuˈbaɾtə ˌnɔbəˈmen} \\
9.11 & My best wishes to you and your family. & Els meus millors desitjos per a vostè i la seva família. & \ipa{əls ˈmɛws miˈʎos dəˈzidʒus pəɾ ə busˈtɛ i lə ˈsebə fəˈmiliə} \\
9.12 & Best wishes to your wife. & Els millors desitjos per a la seva esposa. & \ipa{əls miˈʎos dəˈzidʒus pəɾ ə lə ˈsebə əsˈpɔzə} \\
9.121 & Best wishes to your husband. & Els millors desitjos per al seu marit. & \ipa{əls miˈʎos dəˈzidʒus pəɾ əl ˈsɛw məˈɾit} \\
9.13 & Good luck. & Bona sort. & \ipa{ˈbɔnə ˈsɔɾt} \\
9.14 & Good health. & Estigui bonet. & \ipa{əsˈtiɡi buˈnɛt} \\
9.15 & Goodbye & Adéu. & \ipa{əˈdɛw} \\
9.16 & I wish you an enjoyable afternoon. & Li desitjo una bona tarda. (12–18 h) & \ipa{li dəˈzidʒu ˈunə ˈbɔnə ˈtaɾdə} \\
9.161 & I wish you an enjoyable evening. & Li desitjo un bon vespre. & \ipa{li dəˈzidʒu un ˈbɔn ˈbɛspɾə} \\
9.162 & I wish you an enjoyable weekend. & Li desitjo un bon cap de setmana. & \ipa{li dəˈzidʒu un ˈbɔn ˈkab də səmˈmanə} \\
9.17 & Enjoy your holiday. & Frueixi de les vacances. & \ipa{fɾuˈɛʃi də ləs bəˈkansəs} \\
9.18 & Happy New Year. & Bon any nou. & \ipa{ˈbɔn ˈaɲ ˈnɔw} \\
\rowcolor{band} 10.0 & \textbf{Country} & \multicolumn{2}{l}{\textbf{País}} \\*
10.1 & Albania & Albània & \ipa{əlˈbaniə} \\
10.2 & Andorra & Andorra & \ipa{ənˈdɔrə} \\
10.3 & Austria & Àustria & \ipa{ˈawstɾiə} \\
10.4 & Belgium & Bèlgica & \ipa{ˈbɛlʒikə} \\
10.5 & Belorussia & Bielorússia & \ipa{biəluˈrusiə} \\
10.6 & Bosnia \& Herzegovina & Bòsnia i Herzegovina & \ipa{ˈbɔzniə i əɾzəɡuˈbinə} \\
10.7 & Bulgaria & Bulgària & \ipa{bulˈɡaɾiə} \\
10.8 & Catalonia & Catalunya & \ipa{kətəˈluɲə} \\
10.9 & Croatia & Croàcia & \ipa{kɾuˈasiə} \\
10.11 & Czech Republic & República Txeca & \ipa{rəˈpublikə ˈtʃɛkə} \\
10.12 & Denmark & Dinamarca & \ipa{dinəˈmaɾkə} \\
10.13 & Eire & Irlanda & \ipa{iɾˈlandə} \\
10.14 & England & Anglaterra & \ipa{əŋɡləˈtɛrə} \\
10.15 & Estonia & Estònia & \ipa{əsˈtɔniə} \\
10.16 & Finland & Finlàndia & \ipa{finˈlandiə} \\
10.17 & France & França & \ipa{ˈfɾansə} \\
10.18 & Germany & Alemanya & \ipa{ələˈmaɲə} \\
10.19 & Gibraltar & Gibraltar & \ipa{ʒibɾəlˈta} \\
10.21 & Greece & Grècia & \ipa{ˈɡɾɛsiə} \\
10.22 & Holland & Holanda & \ipa{uˈlandə} \\
10.23 & Hungary & Hongria & \ipa{uŋˈɡɾiə} \\
10.24 & Iceland & Islàndia & \ipa{izˈlandiə} \\
10.25 & Northern Ireland & Irlanda del Nord & \ipa{iɾˈlandə dəl ˈnɔɾt} \\
10.26 & Isle of Man & Illa de Man & \ipa{ˈiʎə də ˈman} \\
10.27 & Italy & Itàlia & \ipa{iˈtaliə} \\
10.28 & Latvia & Letònia & \ipa{ləˈtɔniə} \\
10.29 & Liechtenstein & Liechtenstein & \ipa{ˈlɪçtn̩ʃtaɪn} \\
10.31 & Lithuania & Lituània & \ipa{lituˈaniə} \\
10.32 & Luxembourg & Luxemburg & \ipa{luksəmˈbuɾk} \\
10.33 & Macedonia & Macedònia & \ipa{məsəˈdɔniə} \\
10.34 & Moldavia & Moldàvia & \ipa{mulˈdabiə} \\
10.35 & Netherlands & Països Baixos & \ipa{pəˈizus ˈbaʃus} \\
10.36 & Norway & Noruega & \ipa{nuˈɾwɛɡə} \\
10.37 & Poland & Polònia & \ipa{puˈlɔniə} \\
10.38 & Portugal & Portugal & \ipa{puɾtuˈɡal} \\
10.39 & Romania & Romania & \ipa{ruməˈniə} \\
10.41 & Russian Federation & Federació Russa & \ipa{fədəɾəsiˈo ˈrusə} \\
10.42 & Scotland & Escòcia & \ipa{əsˈkɔsiə} \\
10.43 & Slovakia & Eslovàquia & \ipa{əzluˈbakiə} \\
10.44 & Slovenia & Eslovènia & \ipa{əzluˈbɛniə} \\
10.45 & Spain & Espanya & \ipa{əsˈpaɲə} \\
10.46 & Svalbard & Svalbard & \ipa{ˈsvɑːlbɑr} \\
10.47 & Sweden & Suècia & \ipa{suˈɛsiə} \\
10.48 & Switzerland & Suïssa & \ipa{suˈisə} \\
10.49 & Turkey & Turquia & \ipa{tuɾˈkiə} \\
10.50 & Ukraine & Ucraïna & \ipa{ukɾəˈinə} \\
10.51 & Wales & Gales & \ipa{ˈɡaləs} \\
\end{longtable}
\normalsize
\section*{Hinweise zur Übertragung und Aussprache}
\fontsize{8.5}{10}\selectfont\raggedright
\textbf{Vorlage und Umfang.} Grundlage ist „Catalan v 1 0.doc“; Credit: Jaume, EA3FRB.
Alle 144 Phrasen und Ländernamen sind enthalten, einschließlich Andorra und Catalunya.
Die ursprünglichen Referenzen, Nummernsprünge und der fehlende Abschnitt 7 bleiben erhalten.

\smallskip
\textbf{Korrekturen.} „Bonic dia“ in 1.31 wurde gemäß „Good day“ zu „Bon dia“; „llicencia“ zu „llicència“
(2.2/2.3), „es“ zu „és“ (2.5/2.51/2.52/3.1), „radio“ zu „ràdio“ (3.1), „son“ zu „són“ (4.2/4.3).
„tarja/targes“ wurde bei den QSL-Karten zu „targeta/targetes“ (8.1/8.4/8.9/8.11).
Die Zahl 1224 in 8.12 wurde entsprechend der englischen Spalte zu 1234 berichtigt.
„tant agradable“ wurde zu „tan agradable“ (9.2), „per a el“ zu „per al“ (9.121),
„Bielorrusia“ zu „Bielorússia“ (10.5), „Bosnia“ zu „Bòsnia“ (10.6).
Sprachbezeichnung català und modes digitals stehen klein; Leerzeichen und doppelte Punkte wurden vereinheitlicht.

\smallskip
\textbf{Quellenbesonderheiten.} „Estigui bonet“ (9.14) ist eine ältere Abschiedsformel mit einem
Gesundheitswunsch und bleibt erhalten. Die knappen Formulierungen in 8.1 und 8.11 wurden nicht zu vollständigen
Sätzen ergänzt. „l’interior“ (2.8) und „muntanya“ (2.82) entsprechen nicht wortwörtlich der englischen Spalte.
Historische Ländernamen und Webadressen bleiben erhalten. Zeitangaben in Klammern (1.32/9.16) und die
Erklärung „associació“ (8.1) werden in der IPA-Spalte nicht mitgesprochen.

\smallskip
\textbf{Aussprachebasis.} Die manuelle IPA-Spalte ist eine breite Aussprachehilfe für Zentralkatalanisch
im Raum Barcelona. Unbetontes a/e wird überwiegend als \ipa{ə}, o als \ipa{u} dargestellt; betonte offene
\ipa{ɛ, ɔ} unterscheiden sich von \ipa{e, o}. Valencia, Balearen und andere katalanische Regionen können anders klingen.
\ipa{ˈ} markiert Wortbetonung, \ipa{ˌ} ausgewählte Nebenbetonungen, etwa bei Adverbien auf -ment.
Satzmelodie und individuelle Sprechermerkmale sind nicht kodiert.

\smallskip
\textbf{Konsonanten und Reduktionen.} \ipa{ʎ} ist ll, \ipa{ɲ} ny, \ipa{ʃ} der Laut von ix beziehungsweise vielen x;
\ipa{ʒ} steht für j/g vor e/i, \ipa{dʒ} für tj/tg. \ipa{ɾ} ist ein einfacher r-Schlag, \ipa{r} ein gerolltes r.
Schriftliches v wird hier wie b gelesen. Zur Übersicht stehen überwiegend \ipa{b, d, ɡ}; im Satz werden sie häufig
als weichere Approximanten \ipa{β, ð, ɣ} gesprochen. Endverhärtung und einige Ausfälle sind berücksichtigt:
„content“, „moment“, „vint“ und viele Infinitive ohne auslautenden Verschlusslaut beziehungsweise r.
„setmana“ steht mit \ipa{mː}. Weitere Wortverbindungen, Vokalausfälle und Assimilationen sind nicht vollständig dargestellt.

\smallskip
\textbf{Namen und Kürzel.} Sittingbourne, Dover, Icom und Loop stehen mit englischer Aussprache;
Hellschreiber und Liechtenstein deutsch, bureau französisch, Svalbard norwegisch.
Andere Namen sind überwiegend katalanisch gelesen, auch Pere und Gales.
Rufzeichen und Kürzel verwenden katalanische Buchstabennamen: Q „cu“, S „essa“, C „ce“ (\ipa{se}),
G „ge“, J „jota“, H „hac“; W „ve doble“, V „ve baixa“, Y „i grega“.
WACRAL und FISTS werden als Wörter gelesen. Das e in eQSL/eHam ist \ipa{i}; Ham verwendet eine englische Variante,
Webpunkte stehen als „punt“, com und net als Wörter. Namens- und Kürzelvarianten sind möglich.

\smallskip
\textbf{Zahlen und Platzhalter.} 1960 = „mil nou-cents seixanta“, 43 = „quaranta-tres“,
706 = „set-cents sis“, MT63 mit „seixanta-tres“, 599 = „cinc nou nou“,
1234 = „mil dos-cents trenta-quatre“. 3.558 MHz wird im QSO-Kontext als „tres coma cinc cinc vuit megahertz“
gelesen; 80 m, 50 km, 20 cm und 8 °C mit ausgeschriebenen Einheiten.
\texttt{<\#\#\#\#\#>} und offene Adressangaben erscheinen als Auslassung.

\smallskip
\textbf{Fachhinweise.} Joan F. Carbonell und Joaquim Llisterri: „Catalan“, Journal of the International
Phonetic Association 22(1--2), 1992, S. 53--56, DOI 10.1017/S0025100300004618.
J. Llisterri: „Descripció fonètica i fonològica del català oriental central: consonants“, joaquimllisterri.cat.
Optimot: „Tarja, targetó o targeta?“; RodaMots: „estigui bo“ (zur älteren Abschiedsformel).
\end{document}
